\chapter{Quy ước nghiệp vụ}

Đề bài không nói rõ một số điểm, nên trước khi cài đặt đã chốt các quy ước dưới đây. Mọi đoạn mã và kiểm thử đều theo đúng các quy ước này.

\section{Quy ước chung}
\begin{itemize}
  \item \textbf{Mượn cùng lúc} nghĩa là các thiết bị nằm trong \emph{cùng một phiếu mượn}. Phiếu có $k$ thiết bị thì cộng $+1$ cho mọi cặp trong $k$ thiết bị (tức $k(k-1)/2$ cạnh).
  \item \textbf{Quá 24 giờ chưa nhận:} đẩy xuống cuối hàng và tăng \code{missCount}; khi \code{missCount} đạt 2 thì loại hẳn khỏi hàng. Cách này thoả cả ``đẩy xuống cuối'' lẫn ``tự động loại bỏ'' của đề.
  \item \textbf{Đơn vị được giữ cho hàng đợi} không tính vào \code{available}. Bất biến luôn đúng:
  \[\code{available} + \code{borrowed} + \code{onHold} + \code{maintenance} = \code{total}.\]
  \item \textbf{Thời gian} dùng đồng hồ giả lập \code{Clock} (đơn vị giờ), có lệnh tua thời gian để trình bày được logic 24 giờ và quá hạn. Lõi không gọi \code{time()}.
  \item \textbf{Tiền phạt} 36.000 đồng mỗi ngày quá hạn, mỗi thiết bị tính riêng (phiếu 2 thiết bị trễ 3 ngày: $2\times 3\times 36.000 = 216.000$ đồng). Trễ chưa tròn ngày vẫn tính 1 ngày; lưu bằng \code{long long}, không dùng số thực.
  \item \textbf{Không xoá cứng} thiết bị còn phiếu mượn hoặc còn người trong hàng đợi; chỉ được chuyển sang Bảo trì.
\end{itemize}

\section{Các quy ước M1--M19}
\begin{longtable}{@{}l>{\raggedright\arraybackslash}p{13.8cm}@{}}
\caption{Quy ước nghiệp vụ M1--M19}\label{tab:conventions}\\
\toprule
Mã & Nội dung \\
\midrule
\endfirsthead
\toprule
Mã & Nội dung (tiếp) \\
\midrule
\endhead
\bottomrule
\endfoot
\multicolumn{2}{@{}l}{\textit{Thiết bị và trạng thái}}\\
M1 & Trạng thái suy ra từ số lượng: \code{maintenance == total} $\Rightarrow$ Bảo trì; \code{available > 0} $\Rightarrow$ Sẵn sàng; còn lại $\Rightarrow$ Đang mượn. \code{onHold} chỉ là con số, không phải trạng thái. \\
M2 & Bảo trì theo số đơn vị: \code{sendToMaintenance(id, qty)} chỉ lấy từ \code{available}; \code{finishMaintenance(id, qty)} trả về \code{available} rồi phục vụ hàng đợi. \\
M3 & Sửa \code{total} chỉ khi $\code{newTotal} \ge \code{borrowed}+\code{onHold}+\code{maintenance}$; tăng \code{total} thì phục vụ hàng đợi. \\
M4 & Xoá thiết bị (khi đủ điều kiện) xoá cả đỉnh và các cạnh trong đồ thị; phiếu lịch sử giữ nguyên mã, hiển thị ``đã xoá''. \\
M5 & Mã và tên không chứa dấu phẩy, dấu \code{|} hay xuống dòng (lõi kiểm tra) để CSV không cần ngoặc kép. \\
\midrule
\multicolumn{2}{@{}l}{\textit{Mượn / trả}}\\
M6 & Phiếu nguyên tử: phiếu nhiều thiết bị mà có món hết hàng thì thất bại cả phiếu, thông điệp liệt kê món hết; không cập nhật đồ thị. \\
M7 & Mỗi thiết bị tối đa 1 đơn vị trong một phiếu; không trùng mã trong phiếu. \\
M8 & Trả cả phiếu một lần, không trả từng món. \\
M9 & Số ngày mượn hợp lệ 1--30; không giới hạn số phiếu của một sinh viên. \\
M10 & Không có bảng sinh viên riêng; MSSV và họ tên nằm trong phiếu và trong hàng đợi. \\
\midrule
\multicolumn{2}{@{}l}{\textit{Hàng đợi}}\\
M11 & Giữ hàng nhiều người: mỗi đơn vị rảnh được giữ cho người \emph{chưa được giữ} đứng đầu tiếp theo; những người đang giữ luôn là một đoạn đầu của hàng. \\
M12 & Không cho đặt trước khi \code{available > 0} (phải mượn trực tiếp); một sinh viên không đứng hai lần trong cùng một hàng. \\
M13 & Nhận thiết bị (claim) tạo phiếu 1 thiết bị; số ngày nhập lúc nhận (mặc định 7, hợp lệ 1--30). Phiếu 1 món không cộng cạnh đồ thị. \\
M14 & Người bị đẩy xuống cuối mà hàng chỉ còn mình họ thì được giữ lại 24 giờ mới; lỡ lần thứ 2 thì bị loại, đơn vị trở về \code{available}. \\
M15 & Hết hạn dây chuyền: khi tua nhiều giờ, người kế nhận quyền giữ tính từ \code{holdUntil} cũ của người trước (không phải từ ``bây giờ''); lặp đến khi không còn ai hết hạn $\le$ giờ hiện tại. \\
M16 & Có thao tác huỷ đăng ký; nếu người huỷ đang giữ thì chuyển quyền cho người kế. \\
\midrule
\multicolumn{2}{@{}l}{\textit{Đồ thị và khác}}\\
M17 & Top 3: trọng số giảm dần, hoà thì mã tăng dần; ít hơn 3 láng giềng thì trả ít hơn. Top 3 bỏ qua thiết bị Bảo trì; thiết bị hết hàng vẫn gợi ý kèm cờ hết hàng. Gợi ý lúc lập phiếu (\code{suggestForBasket}) cộng trọng số của cùng một thiết bị qua các món trong giỏ, loại món đã có trong giỏ, tối đa 3, tính trên đồ thị \emph{trước khi} cộng phiếu đang tạo. \\
M18 & \code{Clock} đếm giờ kể từ 01/01/2026 00:00. Giờ hiện tại và bộ đếm mã phiếu lưu trong \code{meta.csv}. \\
M19 & Phiếu mở quá hạn hiển thị số ngày quá hạn và phạt tạm tính theo đồng hồ hiện tại; phạt chốt lúc trả. Hạn trả $=$ giờ mượn $+$ số ngày $\times 24$ giờ (không làm tròn tới cuối ngày); quá hạn khi giờ hiện tại $>$ hạn (đúng bằng hạn chưa quá hạn); số ngày quá hạn làm tròn lên. \\
\end{longtable}

\section{Giải thích một số lựa chọn}
\textbf{Vì sao dùng \code{missCount} thay vì chỉ ``đẩy xuống cuối''?} Nếu chỉ đẩy xuống cuối thì một người không bao giờ đến nhận sẽ quay vòng mãi và chiếm chỗ của người khác; nếu loại ngay lần đầu thì không có ``đẩy xuống cuối'' như đề nói. Cho phép lỡ một lần rồi loại ở lần thứ hai là cách dung hoà.

\textbf{Vì sao tách \code{onHold} khỏi \code{available}?} Đơn vị đang được giữ hộ một người thì người khác mượn trực tiếp không được lấy mất. Nếu để lẫn trong \code{available} thì phải tra hàng đợi mỗi lần mượn; tách riêng thì kiểm tra tồn kho chỉ là một phép so sánh $O(1)$, và bất biến tổng bốn số $= \code{total}$ giúp phát hiện lỗi sớm.
